\section{Related Work}
\label{sec:related_work}

The technical foundations of Chatnalyxer span five established research domains: information extraction from informal text, temporal reasoning in conversational settings, multi-modal AI (ASR and OCR), retrieval-augmented generation, and privacy-preserving architectures. This section surveys the relevant literature within each domain, identifies the limitations that persist in current approaches, and motivates the specific design decisions made in Chatnalyxer.

% ─────────────────────────────────────────────────────────
\subsection{Information Extraction and Named Entity Recognition in Informal Text}

Information Extraction (IE) has a decades-long history, but the application of IE techniques to informal, conversational text remains an open challenge. Early systems relied on handcrafted grammars and gazetteers \cite{chiticariu2013}, which achieve high precision on narrow, well-formatted domains but generalize poorly to colloquial language. The transition to statistical sequence labeling—particularly Bi-LSTM-CRF architectures \cite{lample2016}—substantially improved recall on formal corpora; however, instant messaging text violates the lexical and structural assumptions of these models almost universally: utterances are short, punctuation is absent or non-standard, and entities are frequently implied by context rather than named explicitly.

Transformer-based encoders, beginning with BERT \cite{bert2019} and its successors such as RoBERTa \cite{liu2019roberta}, partially mitigate these limitations by capturing long-range bidirectional context. Studies on social media NER benchmarks such as WNUT-17 \cite{derczynski2017} have shown that transformer models substantially outperform BiLSTM-CRF approaches on noisy user-generated text. Nevertheless, even highly competitive NER pipelines face a fundamental limitation in this setting: they operate on fixed entity type taxonomies and require domain-specific training data to generalize. Constructing such datasets for the open-ended diversity of group chat conversations—where tasks may be phrased as questions, imperatives, or offhand references—is prohibitively expensive.

Recent work has framed extraction as a generative task to avoid this dependency. DEGREE \cite{hsu2022degree} formulates event extraction as a conditional text generation problem using template-based prompting, enabling extraction without entity-type-specific training data. Similarly, large language models operating under zero-shot and few-shot prompting regimes (surveyed in \cite{wei2022emergent}) have been shown to extract structured information from free text with competitive accuracy compared to fine-tuned NER models, provided the output format is tightly constrained through the prompt. Chatnalyxer adopts this generative paradigm, constraining the LLM's output via a typed JSON schema rather than relying on a trained extraction model.

% ─────────────────────────────────────────────────────────
\subsection{Temporal Expression Recognition and Normalization}

Temporal information extraction has been formalized through the TimeML annotation framework \cite{pustejovsky2003timeml} and evaluated across successive TempEval shared tasks \cite{uzzaman2013semeval}. Tools such as SUTime \cite{chang2012sutime} and HeidelTime \cite{strotgen2010heideltime} achieve strong results on formal, document-level text by anchoring temporal expressions to a known reference date derived from the document's metadata. However, both systems are rule-based and depend on explicit temporal markers (e.g., ``March 15, 2024'') for reliable resolution.

The dominant form of temporal language in instant messaging is deictic rather than absolute: expressions like ``tomorrow'', ``by EOD'', ``after class'', or ``this weekend'' carry no calendar-specific meaning independent of the speaker's position in time at the moment of utterance. Mazur and Dale \cite{mazur2010wikiwars} demonstrated that deictic temporal resolution requires access to the document creation time as a computational anchor; without it, rule-based systems either produce null outputs or hallucinate plausible dates. Recent work by Xu et al. \cite{xu2021discourse} showed that discourse-aware prompting of LLMs improves deictic resolution in conversational text, but they did not integrate message network transmission timestamps—the most reliable temporal anchor available in a messaging pipeline.

Chatnalyxer addresses this gap directly. By prepending the exact UTC transmission timestamp of each message to the LLM prompt before inference, the system provides a deterministic calendar anchor without any model fine-tuning. The ablation results (Section \ref{sec:ablation_study}) confirm that this single design decision is responsible for over 60\% of the deictic temporal accuracy, which is consistent with the findings of \cite{mazur2010wikiwars} and \cite{xu2021discourse}.

% ─────────────────────────────────────────────────────────
\subsection{Multi-Modal Processing: ASR and OCR in NLP Pipelines}

The integration of Automatic Speech Recognition (ASR) and Optical Character Recognition (OCR) into NLP pipelines has matured considerably. Radford et al. introduced Whisper \cite{radford2022whisper}, an encoder-decoder transformer trained on 680,000 hours of weakly supervised audio data, achieving competitive zero-shot ASR across over 75 languages without per-language fine-tuning. Whisper's architecture demonstrates that large-scale weak supervision over audio can produce transcript quality competitive with human-level baselines on standard benchmarks, making it practical for deployment in heterogeneous audio environments such as group voice note recordings.

In the document understanding domain, Shi et al. \cite{shi2017crnn} and Liao et al. \cite{liao2020mask} established the foundation for end-to-end scene text recognition using convolutional-recurrent hybrid networks and mask-based spotting architectures, respectively. Cloud OCR services—including Microsoft Azure Computer Vision and Google Cloud Vision—operationalize these advances at scale, providing high-accuracy text extraction from images without requiring on-premise model management.

A persistent limitation across both modalities is error propagation: transcription errors (measured as Word Error Rate for ASR and Character Error Rate for OCR) propagate silently into the extraction layer, which receives corrupted text that appears syntactically valid. Li et al. \cite{li2023blip2} and similar multi-modal learning frameworks have explored joint audio-visual representations that could partially address this by disambiguating transcription errors through visual context, but this approach requires substantially more complex model integration than a sequential pipeline. Chatnalyxer adopts the simpler, higher-reliability sequential design while acknowledging the resulting performance ceiling on low-quality audio and image inputs in Section \ref{sec:threats}.

% ─────────────────────────────────────────────────────────
\subsection{Large Language Models for Zero-Shot and Structured Extraction}

The introduction of GPT-3 \cite{brown2020language} established the empirical basis for zero-shot and few-shot task performance through in-context learning, demonstrating that scaling language model parameters enables instruction-following without gradient-based task adaptation. Subsequent instruction-tuned models—including FLAN-T5 \cite{chung2022scaling} and LLaMA-2 \cite{touvron2023llama}—have shown that explicit instruction fine-tuning further improves zero-shot task generalization, including structured extraction from informal text.

Wei et al. \cite{wei2022emergent} characterized chain-of-thought reasoning as an emergent capability of sufficiently large language models, enabling multi-step inference that is particularly relevant for tasks requiring temporal arithmetic (e.g., computing ``next Tuesday'' from a UTC anchor). However, unrestricted generation remains an obstacle for reliable automation: stochastic output can produce malformed JSON, truncated structures, or schema violations. Several approaches have been proposed to address this, including constrained decoding \cite{willard2023efficient} and self-correction loops \cite{madaan2023selfrefine}. Chatnalyxer combines both: the LLM is instructed to generate output conforming to a typed JSON schema, and any validation failure triggers a second inference pass in which the model receives its own malformed output alongside the Python exception trace and is asked to correct it.

% ─────────────────────────────────────────────────────────
\subsection{Retrieval-Augmented Generation for Grounded Conversational AI}

Lewis et al. \cite{lewis2020rag} introduced Retrieval-Augmented Generation (RAG) as a principled method for grounding LLM outputs in external, non-parametric knowledge, demonstrating substantial reductions in factual hallucination compared to purely parametric generation on knowledge-intensive NLP tasks. Subsequent work has extended this framework to conversational settings: \cite{shuster2021retrieval} showed that retrieval grounding reduces entity hallucination in open-domain dialogue by 10--30\%, and \cite{peng2023check} proposed a check-cite-correct pipeline that improves factual consistency in retrieved-augmented responses.

These advances directly motivate the RAG component of Chatnalyxer. Deployed as a conversational assistant over the user's own message and document history, the system faces exactly the hallucination risk that RAG is designed to address: without retrieval grounding, the LLM might answer ``When was that deadline due?'' with a plausible but fabricated date. By retrieving the relevant task record from PostgreSQL and anchoring the generation to that retrieved context, Chatnalyxer ensures that responses are directly traceable to user-authored content.

% ─────────────────────────────────────────────────────────
\subsection{Privacy-Preserving Architectures for Messaging Analytics}

Systems that process private communication data face a heightened threat model: a passive listener that retains message content is a high-value surveillance target regardless of its stated purpose. Dwork and Roth \cite{dwork2014algorithmic} formalized differential privacy as a mathematically rigorous framework for computing aggregate statistics from private data while providing formal per-record privacy properties; however, applying differential privacy to real-time extraction pipelines incurs latency overhead and accuracy degradation that is impractical at the per-message level.

Marcelino et al. \cite{marcelino2021privacy} and related work on privacy-preserving NLP have explored data minimization as an architectural strategy: rather than applying noise mechanisms, systems can achieve strong privacy characteristics by simply not retaining the sensitive input after processing. This architectural approach—ephemeral processing with persistent metadata-only storage—is exactly the model adopted in Chatnalyxer. Raw message content is held only in volatile RAM during the inference window; the database schema enforces storage of only the extracted task name, deadline, and priority score. One-way hashing of group identifiers \cite{boneh2000sok} ensures that even a full database exposure cannot reconstruct which conversation produced a given task record. The privacy model does not provide differential privacy in the strict Dwork-Roth sense, but it substantially limits the primary operational risk of passive message retention.

% ─────────────────────────────────────────────────────────
\subsection{Asynchronous Microservice Architectures for Real-Time NLP}

Real-time NLP pipelines that process high-frequency event streams must decouple ingestion latency from inference latency to remain stable under burst traffic. This requirement is well-established in the stream processing literature: Kulkarni et al. \cite{kulkarni2015twitter} demonstrated with Twitter Heron that synchronous coupling between ingestion and heavy computation stages leads to predictable backpressure failures under Poisson-distributed arrival rates. The HTTP 202 Accepted pattern \cite{fielding2014http} generalizes this decoupling to RESTful microservice architectures, allowing a receiving service to immediately acknowledge a request while processing it asynchronously—bounding the observed ingestion latency to local I/O time rather than downstream compute time.

This engineering principle directly informs the Chatnalyxer architecture. The Node.js edge client—implemented using the Baileys Signal-protocol library—achieves effective E2EE message interception without requiring access to WhatsApp's server-side infrastructure. The immediate HTTP 202 dispatch to the FastAPI core ensures that edge-side latency remains sub-15 ms regardless of LLM API response time, a behavior that the ablation study (Section \ref{sec:ablation_study}) confirms is an operational prerequisite for deployment in active group chat environments.

% ─────────────────────────────────────────────────────────
\subsection{Research Gap}

Table \ref{tab:related_work_comparison} situates Chatnalyxer within the landscape of directly comparable research systems. Examining the prior literature collectively reveals a consistent pattern of partial solutions. IE-focused systems \cite{lample2016,hsu2022degree} address extraction but assume text-only input, operate on static datasets, and require explicit temporal markers for reliable deadline resolution. ASR and OCR pipelines \cite{radford2022whisper,shi2017crnn,liao2020mask} provide robust modality normalization but are not integrated with downstream task extraction. RAG frameworks \cite{lewis2020rag,shuster2021retrieval} improve factual grounding in generation but do not address the passive, ambient extraction problem. To the best of our knowledge, few published systems combine all four capabilities—passive E2EE monitoring, multi-modal normalization, deictic temporal resolution via UTC anchoring, and context-grounded conversational retrieval—within a single, privacy-respecting, real-time pipeline operating over consumer instant messaging infrastructure.

Human-facing productivity tools (calendar applications, enterprise task managers) are deliberately excluded from this comparison: they require explicit user action to create tasks, which is precisely the interaction cost that Chatnalyxer is designed to eliminate. The research gap is therefore not an incremental improvement over any single prior system but an architectural composition that, to the best of our knowledge, has not been previously demonstrated in the published literature.
